% ============================================================
% question-block.tex
% 문제 블록 (번호 배지 + 발문 + 보기/조건 + 선택지 + 답안란)
% 모든 구성을 하나의 minipage로 묶어 단/페이지 단절 금지
%
% 사용:
%   \begin{kfQBlock}{1}{객관식}
%     \kfQStem{다음 글에 대한 설명으로 가장 적절한 것은?}
%     \begin{kfEvidence} ... \end{kfEvidence}    % 선택
%     \begin{kfCondition} ... \end{kfCondition}  % 선택
%     \begin{kfChoices}
%       \kfChoice ...
%     \end{kfChoices}
%     \kfAnswerLines{2}                            % 선택
%   \end{kfQBlock}
% ============================================================

% 무단절 minipage로 묶고, 발문/보기/조건/선택지/답안란을 자유롭게 배치
% 사용자 피드백 #4: [객관식]/[서술형] 등 텍스트 라벨 제거 — 번호 배지만 표시
% #2(유형 라벨)는 호환성 인자로만 받고 출력하지 않음.
\newenvironment{kfQBlock}[2]%
  {\par\smallskip\needspace{6\baselineskip}%
   \noindent\begin{minipage}{\linewidth}%
   \samepage%
   \par\noindent
   \kfQNum{#1}\hspace{4pt}%
   \begin{minipage}[t]{\dimexpr\linewidth-30pt\relax}%
  }%
  {\end{minipage}\end{minipage}\par\smallskip}

% 문제 발문
\newcommand{\kfQStem}[1]{%
  \par\noindent#1\par\vspace{2pt}%
}

% 짧은 소문항 라벨 (1), (2), (3)
\newcounter{kfSubQ}
\newcommand{\kfSubQReset}{\setcounter{kfSubQ}{0}}
\newcommand{\kfSubQ}{%
  \stepcounter{kfSubQ}%
  \par\noindent\textbf{(\arabic{kfSubQ})}~\ignorespaces%
}
